\subsection{A finite ordinary invariant and the remaining plus equations}

\begin{result}{E.148}{Proved}{Only the old-support plus equations remain in the finite ordinary extension route}
Let \(n\ge2\), and retain the hypotheses on \(S\),
the actual arithmetic action, and the five/eleven
Frattini inertia vectors from E.147.
Call the following data admissible at level \(n-1\):
\begin{itemize}
\item a literal global homomorphism in \(Q_{n-1}\),
unramified outside \(S\);
\item specified actual point matches at every old
place other than three, and the ordinary abelian
class of a fixed point \(t_3\) at three;
\item at every finite \(p\notin S\), a lower
abelian class represented by an ordinary
good-reduction point, so both ordinary
valuations are zero.
\end{itemize}
All abelian classes in these data are taken modulo
\(3^{n-1}\). The last requirement is a finite
test at each fixed level, by the bound below.

For such data, there is an admissible extension
at level \(n\), with the same literal predecessor
and the same old point data, if and only if
the quotient \(Q_n/B_n\) has a global lift
unramified outside \(S\) matching the old
point classes away from three.
When this condition holds, the extension can
introduce two or three fresh primes, all
congruent to two modulo nine and outside
any prescribed finite set.
It freezes at each new prime one of the
ordinary integer patterns in (E.147.2).
The new primes may all be taken greater than
\[
                H_{\rm ab}(n)=(3^{2n}+3)^2 .
 \tag{E.148.1}
\]
At every remaining zero-support prime,
the new abelian class still has a
good-reduction ordinary representative.
No actual higher point match at those
zero-support places or at three is required.

The quotient condition itself has only the
following finite obstruction.
Write
\[
 W_n/B_n=C_n^+\oplus C_n^-,
 \qquad C_n^+\simeq\mathbb F_3^{a_n}.
 \]
Choose any global weak quotient lift and,
at each old place, an aligning element
in the finite lower geometric group.
Let
\(\Delta_q^+\in\mathbb F_3^{a_n}\)
be the plus part of the aligned local
point discrepancy, evaluated on Frobenius.
For at least one of the finitely many
aligning-frame choices, there must exist
\(z\in\mathbb F_3^{a_n}\) such that
\[
 \Delta_q^+=k_qz\quad(q\in S,\ q\ne3,\infty),
 \qquad
 k_q=\frac{q^4-1}{3}\pmod3.
 \tag{E.148.2}
\]
This is necessary and sufficient.
Since \(k_5=1\), it can equivalently be
written
\[
                         \Delta_q^+=k_q\Delta_5^+.
 \tag{E.148.3}
\]
The minus-character equations, coefficient-three
ordinary abelian digits, and all zero-support
ordinary abelian conditions impose no additional
obstruction to this one step.

If these finite systems can be passed indefinitely,
the construction gives one continuous full
pro-three section with an infinite-support
three-primary Kummer pair whose local class at
every place comes from an ordinary point.
This implication is conditional: no assertion
is made that the successive systems have a
compatible infinite solution. Nor does a
three-primary construction by itself supply compatibility
with the other coefficient primes in the ordinary pair.
E.151 supplies a separate arithmetic refinement retaining
that full ordinary compatibility, still conditional on an
actual infinite path. Neither result constructs sections
at the other coefficient primes required by the full
pronilpotent question.
\end{result}
\begin{proof}
We first justify the exact form of the quotient
gate. Its kernel is the direct sum
\(C_n^+\oplus C_n^-\) from E.144, with the
trivial and cyclotomic-modulo-three actions.
Its obstruction always vanishes over the
old \(G_{\mathbb Q,S}\).
For the plus part this is E.66 and E.142.
For the minus part it is the Brauer argument:
the old actual points kill all local
invariants away from three, and reciprocity
then kills the invariant at three.
Thus a global weak quotient lift exists
before old local matching is imposed.

The minus first-cohomology localization
at the old primes is surjective:
each coordinate is a valuation class,
independently attained by \(\kappa(q)\).
The plus first cohomology over the old
ramification set has the single scalar
generator
\[
 \bar\eta_4
       =(\chi_3^4-1)/3\pmod3
           \quad\text{in }H^1(G_{\mathbb Q,S},\mathbb F_3).
 \tag{E.148.4}
\]
At \(q\equiv-1\pmod3\), its local first
cohomology is unramified and its Frobenius
value is \(k_q\). Thus, in \(a_n\) copies,
the precise image is \(q\mapsto k_qz\).
The local minus corrections do not affect
the plus component, since the kernel
is this direct sum as a Galois module.
This proves (E.148.2).
All finite projected stabilizers are
retained by testing every aligning
frame in the finite lower group.
Different lifts of one aligning frame
change its discrepancy only by a
kernel coboundary. Thus the criterion
is a finite collection of finite affine
tests, rather than a claim that those
stabilizers vanish.
The formula \(k_5=(5^4-1)/3=208=1\bmod3\)
proves (E.148.3).

Suppose the gate passes and fix a matching
quotient lift. As in E.146, its \(B_n\)
obstruction vanishes by Brauer reciprocity,
so choose a full \(Q_n\)-lift unramified
outside \(S\). Correct the old support
point discrepancies and the valuation
coordinates at three by global \(S\)-unit
Kummer twists.
This leaves the two ordinary abelian
unit digits at three, just as in E.147.
There is no full point condition at three.

The difference from E.147 is at the
zero-support primes: no higher quotient
point condition is needed there.
Set \(N=3^n\) and consider only the
finite set of \(p\notin S\) with
\(p\le H_{\rm ab}(n)\).
By admissibility choose a good-reduction
ordinary point \(t_p\) whose abelianization
agrees with the lower class modulo \(N/3\).
Compare its level-\(N\) abelian class with
the current lift. Both classes are
unramified, and their difference reduces
to zero modulo \(N/3\).
The coefficient sequence
\[
 0\longrightarrow\mu_3\longrightarrow\mu_N
       \longrightarrow\mu_{N/3}\longrightarrow0
 \tag{E.148.5}
\]
therefore supplies a correcting
\(H^1(\mathbb Q_p,\mu_3)^2\)-class.
Here \(H^0(\mathbb Q_p,\mu_{N/3})\)
may be nonzero. We use only existence
of a preimage, not uniqueness.

Every such preimage is unramified.
Indeed inertia acts trivially on these
modules for \(p\ne3\), and the coefficient
inclusion induces an injection
\[
 H^1(I_p,\mu_3)\longrightarrow H^1(I_p,\mu_N):
 \tag{E.148.6}
\]
both are homomorphism groups and the
coefficient map is injective.
The image of the preimage has zero
inertia, so the preimage has zero
inertia too.
Consequently it is a unit Kummer class.
At \(p\equiv-1\pmod3\) this class is zero;
at \(p\equiv1\pmod3\) its two coordinates
are exactly the unit corrections handled
by the cubic Dirichlet fields in E.147.
The possible \(H^0\)-ambiguity in
(E.148.5) is harmless.

Apply that construction with these
abelian unit corrections in place of
the full point discrepancies at the
finite set of zero-support primes.
The quotient/action field remains
unramified there, so the same independent
cubic Dirichlet fields permit every
needed unit prescription.
The two or three spanning peripheral
vectors of (E.147.8)--(E.147.9) correct
all those unit classes and the local-three
unit vector simultaneously.
Choose all new primes above
\(H_{\rm ab}(n)\), with residue two
modulo nine and the required cusp
Frobenius. At each new prime the full
local \(Q_n\)-class is the specified
ordinary point, by E.146, including
the infinity case.
Every old actual support point remains
matched. At the tested zero-support
primes, the abelian class now agrees
with that of the chosen \(t_p\).
The higher nonabelian class there is
not asserted to be a point image.

For completeness, the smaller uniform
bound (E.148.1) follows directly from
twisted Fermat curves.
At any remaining \(p\ne3\), an unramified
abelian class modulo \(N\) has two
unit representatives \(a,b\).
Consider over \(\mathbb F_p\) the curve
\[
                  \bar aX^N+\bar bY^N=Z^N.
 \tag{E.148.7}
\]
It is a smooth, geometrically connected
degree-\(N\) Fermat twist because
\(p\nmid N\) and \(\bar a,\bar b\ne0\).
Its genus and number of geometric
coordinate-boundary points are
\[
          2g=(N-1)(N-2),\qquad b_{\partial}=3N.
 \tag{E.148.8}
\]
Alternatively these follow from the
degree-\(N^2\) Kummer cover of the
tripod, whose three ramification
indices are all \(N\).
The Weil bound
\citep[(1.5.1)--(1.6), pp.275--276]{RootDeligne1974}
gives a rational point
with all coordinates nonzero whenever
\[
 p+1-(N-1)(N-2)\sqrt p-3N>0.
 \]
The sufficient inequality
\(\sqrt p>N^2+3=2g+b_{\partial}+1\)
is exactly (E.148.1).
Such a point realizes the two required
unit power classes by
\[
 t=\bar a(X/Z)^N,\qquad
 1-t=\bar b(Y/Z)^N .
 \]
A lift of \(t\) to \(\mathbb Z_p\)
is a good-reduction ordinary point
with those same local Kummer classes:
taking \(N\)-th powers is an automorphism
on the principal units because \(p\nmid N\).
This proves all the remaining local
abelian conditions.
It also shows that admissibility at
any fixed level requires only finitely
many zero-support tests.

Conversely, any next-level extension
with actual ordinary images at its
new support primes must pass the
quotient gate by E.144.
That argument allows arbitrary new
primes, not only the special primes
selected in our construction.
This proves the exact one-step
equivalence and the closed finite-level
invariant.

We make the good-reduction qualifier
explicit with a local scope check.
At seven the point \(7^3\) has trivial
two-component Kummer class modulo three:
its first coordinate is a cube, and
its complementary unit reduces to one.
But there is no good-reduction point
with that trivial pair, because the
nonzero cubes in \(\mathbb F_7\) are
only \(1,-1\), and no two of them
sum to one. Thus finite ordinary
representability alone cannot be
silently substituted for the
zero-support good-reduction hypothesis.

Finally suppose an infinite sequence
of these extensions has actually
been chosen.
Their literal compatibility gives
a continuous section into
\[
                          F=\varprojlim_n Q_n;
 \]
the lower exponent-three central
series is cofinal among the open
normal subgroups of a finitely
generated pro-three group.
Continuity follows coordinatewise
from the finite continuous sections.
At every prime introduced at stage
\(n\), its external point target
fixes an ordinary integer valuation
divisible by \(3^{n-1}\), and it
remains fixed thereafter.
The union of the newly introduced
primes is infinite, while those
integer valuation coordinates tend
to zero three-adically.
Thus the limiting abelian class is
one three-primary completed Kummer
pair with infinite support, not a
different pair at each level.

At three it is the full Kummer pair
of the fixed \(t_3\).
At an eventual support prime the
fixed point target gives the full
local pair. The finite conjugacy
conditions for its section have
nonempty nested compact sets of
geometric conjugators, so they also
give the full local section match.
At a prime never introduced into
the support, the sets of
good-reduction points realizing
the successive abelian reductions
are nonempty nested closed subsets
of the compact set
\[
 \{t\in\mathbb Z_p:
           t\bmod p\notin\{0,1\}\}.
 \]
Their intersection gives an actual
ordinary local point realizing
the full three-primary pair.
Real three-primary abelianization
is trivial.
This proves the stated conditional
limit implication. It does not
construct the requisite infinite
path through the systems
(E.148.2).
\end{proof}
